\newpage

\section{等差数列}

\begin{definition}
    如果一个数列从第二项起，每一项与它的前一项的差等于同一个常数，这个数列就叫做等差数列，这个常数叫做等差数列的公差，通常用字母 $d$ 表示。

    \begin{tips}
        设等差数列 $\{a_n\}$ 的首项为 $a_1$，公差为 $d$。

        根据等差数列的定义，对于任何 $n \geq 2$，都有 $a_n - a_{n-1} = d$。
        由此，我们可以写下一系列等式：

        $$
            \begin{cases}
                a_n - a_{n-1} = d     \\
                a_{n-1} - a_{n-2} = d \\
                \qquad \vdots         \\
                a_3 - a_2 = d         \\
                a_2 - a_1 = d
            \end{cases}
        $$

        所有式子相加可得：
        $$ (a_n - a_{n-1}) + (a_{n-1} - a_{n-2}) + \dots + (a_3 - a_2) + (a_2 - a_1) =  \underbrace{d + d + \dots + d + d}_{n-1 \text{ 个}} $$
        $$ a_n - a_1 = (n-1)d $$

        将 $-a_1$ 移项至等式右边，即可得到最终的等差数列通项公式：
        $$ a_n = a_1 + (n-1)d $$
    \end{tips}

\end{definition}

\begin{theorem}
    等差数列的通项公式为 $a_n = a_1 + (n-1)d$。

    等差数列的前 $n$ 项和公式为 $S_n = \frac{n(a_1+a_n)}{2} = na_1 + \frac{n(n-1)}{2}d$。
\end{theorem}

\begin{proof}
    \textbf{采用倒序相加法证明：} 设等差数列 $\{a_n\}$ 的前 $n$ 项和为 $S_n$。
    \begin{align}
        S_n & = a_1 + a_2 + a_3 + \cdots + a_{n-1} + a_n                                \\
            & = a_1 + [a_1 + d] + [a_1 + 2d] + \cdots + [a_1 + (n-2)d] + [a_1 + (n-1)d]
    \end{align}

    将次序颠倒，得：
    \begin{align}
        S_n & = a_n + a_{n-1} + a_{n-2} + \cdots + a_2 + a_1               \\
            & = [a_1 + (n-1)d] + [a_1 + (n-2)d] + \cdots + [a_1 + d] + a_1
    \end{align}

    两式相加：
    \begin{align}
        2S_n & = [2a_1 + (n-1)d] + [2a_1 + (n-1)d] + \cdots + [2a_1 + (n-1)d] \\
             & = n[2a_1 + (n-1)d] = n(a_1 + a_n)
    \end{align}

    因此：
    $$S_n = \frac{n(a_1+a_n)}{2} = na_1 + \frac{n(n-1)}{2}d$$
\end{proof}

\subsection{将 $a_{?}$ 和 $S_{?}$ 化为 $a_1$ 和 $d$}

\begin{tips}
    求出数列通项，这个数列的"解析式"就出来了。所以只要求出 $a_1$ 和 $d$，就能得到 $S_n$ 和 $a_n$，能够研究数列（函数）的解析式，这是非常好的！
\end{tips}

\begin{example}
    在等差数列 $\{a_n\}$ 中，已知 $a_3 = 7$，$a_8 = 22$。则该数列的第10项 $a_{10}$ 为 \underline{\qquad}
\end{example}

\v{6em}

\begin{example}
    已知等差数列 $\{a_n\}$ 的前 $n$ 项和为 $S_n$，若 $a_1+a_9 = 10$，则 $S_9$ 的值为（\qquad）
    \begin{tasks}(4)
        \task 40
        \task 45
        \task 50
        \task 90
    \end{tasks}
\end{example}

\v{6em}

\begin{example}
    在等差数列 $\{a_n\}$ 中，已知 $a_5 = 10$ 且前 7 项和 $S_7 = 49$。求该数列的公差 $d$ 和首项 $a_1$。
\end{example}

\v{10em}

\begin{example}
    求数列 $\{a_n\}$ 的前 $n$ 项和 $S_n$，其中 $a_n = 2n - 1$。
\end{example}

\v{10em}

\subsection{$S_n$ 公式的简单运用}

\begin{example}
    已知等差数列 $\{a_n\}$ 中，$a_1 + a_{100} = 202$，求 $S_{100}$。
\end{example}

\v{8em}

\subsection{等差数列的下标和性质}

\begin{theorem}[下标和性质]
    若 $m+n=p+q$，则 $a_m + a_n = a_p + a_q$。

    \textbf{证明：}
    根据等差数列通项公式：
    \begin{align}
        a_m + a_n & = [a_1 + (m-1)d] + [a_1 + (n-1)d] \\
                  & = 2a_1 + (m+n-2)d
    \end{align}

    同理：
    \begin{align}
        a_p + a_q & = [a_1 + (p-1)d] + [a_1 + (q-1)d] \\
                  & = 2a_1 + (p+q-2)d
    \end{align}

    由于 $m+n=p+q$，所以 $a_m + a_n = a_p + a_q$。
\end{theorem}

\begin{example}
    已知 $S_n$ 是等差数列 $\left\{a_n\right\}$ 的前 $n$ 项和，且 $a_7>0, ~ a_6+a_9<0$ ，则 （\qquad）
    \begin{tasks}(4)
        \task 数列 $\left\{a_n\right\}$ 为递减数列
        \task $a_8>0$
        \task $S_n$ 的最大值为 $S_7$
        \task $S_{14}>0$
    \end{tasks}
\end{example}

\begin{tips}
    C 选项在学习下一个知识点后再做。
\end{tips}

\v{12em}

\subsection{等差数列的第三求和公式}

\begin{corollary}[等差数列的第三求和公式]
    根据 $S_n=\frac{(a_1+a_n)}{2}n$，可得：

    $$S_{2n-1}=(2n-1)a_n$$
    或者记为：
    $$S_n = na_{\text{中}}$$
\end{corollary}

\begin{example}[使 $a_n$ 和 $S_n$ 之间相互转化]
    已知 $S_{13}=13$，求 $a_7 = $ \underline{\qquad}.
\end{example}

\v{4em}

\begin{example}
    若两个等差数列 $\left\{a_n\right\}$ 和 $\left\{b_n\right\}$ 的前 $n$ 项和分别是 $S_n$、$T_n$ ，已知 $\frac{S_n}{T_n}=\frac{3 n}{n+2}$ ，则 $\frac{a_7}{b_7}=$ \underline{\qquad}.
\end{example}

\v{6em}

\subsection{等差数列 $S_n$ 和 $a_n$ 的图像关系}

\begin{theorem}[$S_n$ 和 $a_n$ 的图像关系]
    当 $a_n$ 中的 $d<0$ 时，由于 $a_n$ 可以看成一个一次函数，故在 $a_n > 0$ 成立的 $n$ 取最大值时，$S_n$ 取到最大值：
    \begin{center}
        \begin{tikzpicture}[scale=0.6]
            % Axes
            \draw[->] (0,0) -- (8,0) node[right] {$n$};
            \draw[->] (0,-1) -- (0,5) node[above] {$y$};

            % a_n graph (straight line)
            \draw[thick, blue] (0,3) -- (8,-1);

            % Add vertical dashed line at n_0
            \draw[dashed] (5,0) -- (5,0.5);
            \node[below] at (5,0) {$n_0$};

            % When a_n > 0 and a_{n+1} < 0
            \filldraw[blue] (5,0.5) circle (2pt) node[above right] {$a_{n_0} > 0$};
            \filldraw[blue] (7,-0.5) circle (2pt) node[right] {$a_{n_{0+1}} < 0$};
        \end{tikzpicture}
    \end{center}
\end{theorem}

\begin{example}
    （ 2022 春•上饶月考）设等差数列 $\left\{a_n\right\}$ 的前 $n$ 项和为 $S_n$ ，且 $S_{4045}>0, S_{4044}<0$ ，则当 $n=$ \underline{\qquad} 时，$S_n$ 最小.
\end{example}

\v{8em}

\begin{example}
    若数列 $\left\{a_n\right\}$ 为等差数列，$S_n$ 为数列 $\left\{a_n\right\}$ 的前 $n$ 项和，$a_4+a_9>0, ~ S_{11}<0$ ，则 $S_n$ 的最小值为 （\qquad）
    \begin{tasks}(4)
        \task $S_5$
        \task $S_6$
        \task $S_7$
        \task $S_8$
    \end{tasks}
\end{example}

\v{8em}

\begin{example}
    已知 $S_n$ 是等差数列 $\left\{a_n\right\}$ 的前 $n$ 项和，且 $S_6>S_7>S_5$ ，有下列四个命题，其中正确的是 （\qquad）
    \begin{tasks}(2)
        \task $d>0$
        \task $S_{11}>0$
        \task $S_{12}<0$
        \task 数列 $\left\{S_n\right\}$ 中的最大项为 $S_{11}$
    \end{tasks}
\end{example}

\v{8em}

\subsection{等差数列 $S_n$ 的特点}

\begin{theorem}
    $S_n$ 实际上是一个二次函数：$S_n=An^2+Bn+C$，其中，$A=\frac{d}{2}$，$B=a_1-\frac{d}{2}$。

    所以 $S_n = k$ 这样的方程，可能会有两个解。

    \begin{warning}
        $S_n$ 中因为 $n$ 是正整数，所以不一定 $S_n$ 的最大值就在对称轴的位置取。
    \end{warning}

    \begin{center}
        \begin{tikzpicture}[scale=0.8]
            % Axes
            \draw[->] (0,0) -- (6,0) node[right] {$n$};
            \draw[->] (0,-2) -- (0,5) node[above] {$y$};

            % Quadratic function with vertex at (2.5, 4)
            \draw[thick, blue, domain=0:5, smooth, variable=\x]
            plot (\x, {-(\x-2.5)^2 + 4}) node[right] {$S_n = -n^2 + 5n - 2.25$};

            % Mark the vertex
            \filldraw[blue] (2.5,4) circle (2pt) node[above right] {$(2.5, 4)$};

            % Draw dashed lines to show vertex
            \draw[dashed] (2.5,0) -- (2.5,4);
            \draw[dashed] (0,4) -- (2.5,4);

            % Add a horizontal dashed line at y = 2
            \draw[dashed] (0,2) -- (5,2) node[right] {$y = 2$};

            % Label the axes
            \foreach \x in {1,2,3,4,5}
            \draw (\x,0.1) -- (\x,-0.1) node[below] {$\x$};
            \foreach \y in {1,2,3,4}
            \draw (0.1,\y) -- (-0.1,\y) node[left] {$\y$};
        \end{tikzpicture}
    \end{center}
\end{theorem}

\begin{example}
    已知等差数列 $\left\{a_n\right\}$ ，前 $n$ 项和 $S_n=n^2-4 n$ ，则 （\qquad）
    \begin{tasks}(2)
        \task $a_1=-3$
        \task $d=-2$
        \task $a_{10}=15$
        \task $\left\{\frac{S_n}{n}\right\}$ 为公差为 -2 的等差数列
    \end{tasks}
\end{example}

\v{12em}

\begin{example}
    设等差数列 $\left\{a_n\right\}$ 的前 $n$ 项和为 $S_n$ ，公差为 $d, ~ a_1>0, ~ a_6+a_7>0, ~ a_6 \cdot a_7<0$ ，下列结论正确的是 （\qquad）
    \begin{tasks}(4)
        \task $a_6<0, a_7>0$
        \task $d<0$
        \task $S_{13}<0$
        \task 当 $n=7$ 时，$S_n$ 最大
    \end{tasks}
\end{example}

\v{8em}

\begin{example}
    设等差数列 $\left\{a_n\right\}$ 的前 $n$ 项和为 $S_n$ ，公差为 $d$ ，若 $a_1=30, S_{12}=S_{19}$ ，则 （\qquad）
    \begin{tasks}(4)
        \task $d=-2$
        \task $S_n \leq S_{15}$
        \task $a_{15}=2$
        \task $S_{30}=0$
    \end{tasks}
\end{example}

\v{10em}

\begin{example}
    已知数列 $\left\{a_n\right\}$ 的前 $n$ 项和为 $S_n=\frac{1}{2} n^2-\frac{13}{2} n+6$ ，则下列说法正确的是 （\qquad）
    \begin{tasks}(2)
        \task $a_n=n-7$
        \task $\frac{1}{a_2 a_3}+\frac{1}{a_3 a_4}+\frac{1}{a_4 a_5}+\frac{1}{a_5 a_6}=\frac{4}{5}$
        \task 使 $S_n>0$ 的最小正整数 $n$ 为 13
        \task $\frac{S_n}{n}$ 的最小值为 -3
    \end{tasks}
\end{example}

\v{10em}

\subsection{等差中项}

\begin{definition}
    在三个数 $a$, $b$, $c$ 之间，如果 $b = \frac{a+c}{2}$，则称 $b$ 是 $a$ 和 $c$ 的等差中项。
\end{definition}

\begin{theorem}
    若 $a$, $b$, $c$ 成等差数列，则 $b = \frac{a+c}{2}$。

    若 $a$, $b_1$, $b_2$, $\ldots$, $b_n$, $c$ 成等差数列，则有：
    \begin{align}
        b_1 + b_2 + \ldots + b_n = \frac{n(a+c)}{2}
    \end{align}
\end{theorem}

\begin{example}
    \begin{enumerate}
        \item 已知数列 $\left\{a_n\right\}$ 的前 $n$ 项和为 $S_n=n^2+n-1,\left(n \in \mathbf{N}^*\right)$ ，求该数列的通项公式 $a_n$ ．
        \item 已知数列 $\left\{a_n\right\}$ 满足 $2 a_{n+1}=a_n+a_{n+2},\left(n \in \mathbf{N}^*\right)$ ，且 $a_1+a_3+a_5=96, a_2+a_4+a_6=90$ ，求 $a_{20}$ ．
    \end{enumerate}
\end{example}

\v{10em}

\begin{example}
    等差数列 $\left\{a_n\right\}$ 的前 $n$ 项和记为 $S_n$ ，已知 $S_3=-15$ ，且 $a_1, a_3,-a_4$ 成等差数列．
    \begin{enumerate}
        \item 求数列 $\left\{a_n\right\}$ 的通项公式；
        \item 求数列 $\left\{a_n\right\}$ 的前 $n$ 项和 $S_n$
    \end{enumerate}
\end{example}

\v{12em}

\subsection{等差数列的公共项}

\begin{theorem}[等差数列的公共项]
    设数列 $a_n=an+b$ 和 $b_n=cn+d$ 分别为等差数列，它们的公共项组成一个新的数列。
    \begin{enumerate}
        \item 公共项数列存在当且仅当 $a, c$ 的最大公因数（记为 $(a, c)$ ）整除 $d-b$ ，更简洁一点就是 $(a, c) \mid d-b$ 。这个可以通过 $a x+b=c y+d$ 有正整数解，然后由 Bezout 定理相关结论得到。
        \item 如果存在，公共项数列 $\left\{C_n\right\}$ 的通项是
    \end{enumerate}
    $$
        C_n=\frac{a c}{(a, c)} n+C_1-\frac{a c}{(a, c)}
    $$
\end{theorem}

\begin{tips}
    当然了，最简单的应该还是直接列举出来找规律。
\end{tips}

\begin{example}
    （2020·海南·15）将数列 $\{2 n-1\}$ 与 $\{3 n-2\}$ 的公共项从小到大排列得到数列 $\left\{a_n\right\}$ ，则 $\left\{a_n\right\}$ 的前 $n$ 项和为 \underline{$\quad$} ．
\end{example}

\vspace{6em}

\subsection{证明一个数列是等差数列}

\begin{theorem}
    \begin{itemize}
        \item 需要严格证明后一项和前一项的差为定值，这才是等差数列的基础概念。
        \item 或者证明等差中项的式子 $2 a_{n+1} = a_{n} + a_{n+2}$
    \end{itemize}
\end{theorem}

\begin{example}
    设 $S_n$ 为数列 $\left\{a_n\right\}$ 的前 $n$ 项和，$S_n=2 n^2-30 n$ ．
    \begin{enumerate}
        \item 求 $a_n$ ；

        \item 证明 $\left\{a_n\right\}$ 是等差数列．
    \end{enumerate}
\end{example}

\v{12em}

\begin{example}
    （2022－甲卷（文））记 $S_n$ 为数列 $\left\{a_n\right\}$ 的前 $n$ 项和．已知 $\frac{2 S_n}{n}+n=2 a_n+1$ ．
    \begin{enumerate}
        \item 证明：$\left\{a_n\right\}$ 是等差数列；
    \end{enumerate}
\end{example}

\v{10em}

\begin{example}（2016·天津·理 18 题）
    已知 $\left\{a_n\right\}$ 是各项均为正数的等差数列，公差为 $d$ ．对任意的 $n \in \mathbf{N}^*, b_n$ 是 $a_n$ 和 $a_{n+1}$ 的等比中项．
    \begin{enumerate}
        \item （I）设 $c_n=b_{n+1}^2-b_n^2, n \in \mathbf{N}^*$ ，求证：数列 $\left\{c_n\right\}$ 是等差数列；
    \end{enumerate}
\end{example}

\v{10em}



